% SPDX-License-Identifier: CC-BY-4.0
\documentclass[11pt]{amsart}
\usepackage[T1]{fontenc}
\usepackage{lmodern}
\usepackage[margin=1in]{geometry}
\usepackage{amsmath,amssymb}
\usepackage[expansion=false]{microtype}
\usepackage[hidelinks]{hyperref}
\newcommand*{\AuthorName}{Felix Dongting Fei}
\newcommand*{\AuthorAlternativeName}{Dongting Fei}
\newcommand*{\AuthorAffiliation}{Independent Researcher}
\newcommand*{\AuthorEmail}{felix.d.fei@gmail.com}


\newtheorem{theorem}{Theorem}[section]
\newtheorem{lemma}[theorem]{Lemma}
\newtheorem{proposition}[theorem]{Proposition}
\newtheorem{corollary}[theorem]{Corollary}
\theoremstyle{remark}
\newtheorem{remark}[theorem]{Remark}
\numberwithin{equation}{section}
\newcommand{\qbinom}[2]{\genfrac{[}{]}{0pt}{}{#1}{#2}_q}
\newcommand{\NN}{\mathbb Z_{\ge0}}
\newcommand{\coeff}[1]{[q^{#1}]}

\newcommand{\PaperTitle}{Uniform First Differences of Gaussian Polynomials and a Proof of the Stanley\textendash{}Zanello Unimodality Conjecture}
\title[First differences of Gaussian polynomials]{\PaperTitle}
\author{\AuthorName}
\thanks{Alternative name: \AuthorAlternativeName. ORCID: \href{https://orcid.org/0009-0000-8891-7101}{0009-0000-8891-7101}. Licensed under \href{https://creativecommons.org/licenses/by/4.0/}{CC BY 4.0}.}
\ifx\AuthorAffiliation\empty\else\address{\AuthorAffiliation}\fi
\ifx\AuthorEmail\empty\else\email{\AuthorEmail}\fi
\date{26 November 2026}
\subjclass[2020]{05A17, 05A16}
\keywords{Gaussian polynomial, integer partition, unimodality, asymptotic enumeration}
\hypersetup{pdftitle={\PaperTitle},pdfauthor={\AuthorName}}

\begin{document}
\begin{abstract}
We prove the asymptotic unimodality conjecture of Stanley and Zanello for
differences of Gaussian polynomials whose lower indices differ by two.
For every fixed \(k\ge5\), all admissible centered differences are
coefficientwise nonnegative and unimodal once the upper parameter is
sufficiently large, apart from the exception stated in their conjecture.
The proof combines an injection for the edge coefficients with a uniform
estimate for first differences in the remaining range. A two-term
restricted-partition estimate and averaging over subsets of a fixed
cardinality give a cardinal-spline approximation with error
\(O_k(\ell^{k-4})\), where \(\ell\) is the width of the height-\(k\) rectangle.
From degree \(\ell\) to \(\lfloor k\ell/2\rfloor\), its first differences
have a lower bound of order \(\ell^{k-3}\), while those of the shifted
term are uniformly \(O_k(\ell^{k-4})\).
A complete Lean~4 formalization proves Conjecture~1.1(iv) in its original
integer parameters.
\end{abstract}
\maketitle
\enlargethispage{2pt}

\section{Introduction and main result}

For \(r,w\in\NN\), write
\begin{equation}\label{eq:G}
G_{r,w}(q)=\qbinom{r+w}{r}
=\prod_{j=1}^{r}\frac{1-q^{w+j}}{1-q^j}.
\end{equation}
This is the generating polynomial for partitions contained in an
\(r\)-by-\(w\) rectangle. Its coefficients are nonnegative, and it is
symmetric about degree \(rw/2\). The empty product, and also \(G_{r,0}\),
is \(1\). Coefficients outside the support of a polynomial are taken to
be zero.

Reiner and Stanton~\cite{RS98} studied unimodality of differences of
specialized Schur functions, including centered differences of Gaussian
polynomials. Stanley and Zanello~\cite{SZ20} proposed a broader
conjecture. For fixed \(k\ge5\), their Conjecture~1.1(iv) asserts that
\begin{equation}\label{eq:original}
f(k,m,b;q)=\qbinom{m}{k}
-q^{\,k(m-b)/2+b-2k+2}\qbinom{b}{k-2}
\end{equation}
has nonnegative, unimodal coefficients for all sufficiently large integers
\(m\ge k\),
uniformly over integers \(b\) such that
\begin{equation}\label{eq:admissible-original}
k-2\le b\le\frac{km-4k+4}{k-2},
\qquad k(m-b)\equiv0\pmod2,
\end{equation}
except for
\begin{equation}\label{eq:exception-original}
b=\frac{km-4k+2}{k-2}
\end{equation}
when this is an integer.
They proved the case \(k=5\) and characterized the smaller lower
indices. Koutschan, Uncu, and Wong~\cite{KUW23} subsequently proved the
cases \(k=6,7\), with explicit information on small parameters.

We prove the assertion for every fixed \(k\ge5\). The threshold depends only on \(k\); no explicit value is obtained.

\begin{theorem}\label{thm:main}
For each integer \(k\ge5\), there exists a nonnegative integer
\(\ell_0(k)\) such
that the following holds. If \(\ell\) is an integer with
\(\ell\ge\ell_0(k)\) and \(s,t\in\NN\)
satisfy
\begin{equation}\label{eq:alignment}
(k-2)t+2s=k\ell,\qquad s\ne1,
\end{equation}
then
\begin{equation}\label{eq:F}
F(q)=G_{k,\ell}(q)-q^sG_{k-2,t}(q)
\end{equation}
has nonnegative, symmetric, unimodal coefficients.
If \(\ell,t\in\NN\) and
\((k-2)t+2=k\ell\), then the coefficients of \eqref{eq:F},
with \(s=1\), are not unimodal.
\end{theorem}

The substitutions
\begin{equation}\label{eq:substitution}
\ell=m-k,\qquad t=b-k+2,\qquad
s=\frac{k\ell-(k-2)t}{2}
\end{equation}
identify \eqref{eq:F} with \eqref{eq:original}.
Indeed,
\[
  2s=km-4k+4-(k-2)b=k(m-b)+2(b-2k+2).
\]
The upper bound on \(b\) is equivalent to \(2s\ge0\), and the
parity condition makes \(s\) an integer even when \(m-b<0\).
The lower bound on \(b\) is exactly \(t\ge0\).
Conversely, recover \(m=k+\ell\) and \(b=k-2+t\) from any aligned
triple. The exceptional equality is equivalent to \(2s=2\).
A width threshold \(W(k)\) gives the threshold \(k+W(k)\) in \(m\). Thus Theorem~\ref{thm:main} proves
\cite[Conjecture~1.1(iv)]{SZ20}.
The quantifier on \(m\) distinguishes this asymptotic result from the
earlier Reiner--Stanton conjecture at finite parameters.

The asymptotic behavior of Gaussian coefficients with one parameter
fixed was studied by Stanley and Zanello~\cite{SZ16}.
Their Theorem~2.4 gives the truncated-power coefficient profile
that appears below as a cardinal-spline density, and their
Theorem~2.6 gives the first-difference asymptotic
at a fixed distance from the center when \(k\ge6\) is even.
The subset-cardinality averaging used below is also present in the
proof of that theorem. Here we use a shifted restricted-partition
estimate to obtain a remainder uniform over the coefficient index,
for both parities of \(k\). The resulting lower bound, together with
a lower-degree comparison and an edge injection, proves the
conjecture.

An earlier, different conjecture of Stanley and Zanello, concerning
partitions into a bounded number of distinct parts, was proved by
Alpoge~\cite{Alpoge14}. That proof also combines elementary edge
arguments with an analytic estimate involving sums of uniform random
variables. For the centered Gaussian differences in \eqref{eq:original}, the
comparison requires an estimate uniform in the coefficient index.

For equal rectangle widths, $t=\ell$ and $s=\ell$, the difference
specializes to the family studied by Chen and Jia~\cite{CJ}.
For this family, a recent preprint of Kang~\cite{Kang} reports strict
increase from degree one to the middle for even widths at least eight
and heights at least fifteen. Melczer, Panova, and Pemantle~\cite{MPP}
obtain sharp first-difference asymptotics for a single rectangle when
its two side lengths grow at comparable rates. Our estimates allow the
height to stay fixed and the smaller width to vary across the full
admissible range.

In all estimates, constants carrying a subscript \(k\), or implicit
in \(O_k(\cdot)\), depend only on \(k\). In particular, they do not
depend on \(s,t\), or on a coefficient index in the specified range.
We use \(x_+=\max(x,0)\).

\section{A shifted restricted-partition estimate}

For \(r\ge3\), define
\begin{equation}\label{eq:restricted}
H_r(q)=\prod_{j=2}^{r}(1-q^j)^{-1},
\qquad a_r(v)=\coeff{v}H_r(q),
\end{equation}
with \(a_r(v)=0\) for \(v<0\).
General formulas for the polynomial part of a restricted partition
function are classical. We refer to Beck, Gessel, and
Komatsu~\cite{BGK}. Here we need the top two terms together with an
error estimate valid across the break at zero.
For nonnegative \(v\), the number \(a_r(v)\) counts partitions of \(v\)
with parts in \(\{2,\ldots,r\}\).

\begin{lemma}\label{lem:restricted}
Fix \(k\ge5\), and put
\[
d=k-2,\qquad S=\sum_{j=2}^{k}j=\frac{k(k+1)}2-1.
\]
Uniformly over integer \(v\),
\begin{equation}\label{eq:restricted-estimate}
a_k(v)=\frac{(v+S/2)_+^d}{k!\,d!}
+O_k\bigl((1+v_+)^{d-2}\bigr).
\end{equation}
\end{lemma}

\begin{proof}
The rational function \(H_k(q)\) has poles only at roots of unity.
Its pole at \(q=1\) has order \(k-1=d+1\).
Writing \(z=1-q\), we obtain
\[
1-q^j=jz\left(1-\frac{j-1}{2}z+O_j(z^2)\right)
\]
and hence
\begin{equation}\label{eq:root-one}
H_k(q)=\frac{z^{-(d+1)}}{k!}
\left(1+\frac{k(k-1)}4z+O_k(z^2)\right).
\end{equation}
Set \(A=k(k-1)/4\). The highest two terms of the principal part in
\eqref{eq:root-one} contribute
\begin{equation}\label{eq:root-one-coeff}
\frac1{k!}
\left(\binom{v+d}{d}
+A\binom{v+d-1}{d-1}\right)
\end{equation}
to the coefficient of \(q^v\), for \(v\ge0\).
The denominator has larger degree than the numerator, so the partial
fraction decomposition has no polynomial part. All remaining terms at one contribute a polynomial
of degree at most \(d-2\).
The coefficient of \(v^{d-1}\) inside the parentheses in
\eqref{eq:root-one-coeff} is
\[
\frac{(d+1)/2+A}{(d-1)!}
=\frac{S/2}{(d-1)!},
\]
since
\[
\frac{d+1}{2}+A
=\frac{k-1}{2}+\frac{k(k-1)}4
=\frac{k(k+1)-2}{4}.
\]
Therefore the contribution of the pole at \(1\) is
\begin{equation}\label{eq:root-one-shift}
\frac{(v+S/2)^d}{k!\,d!}
+O_k((v+1)^{d-2}).
\end{equation}

For a root of unity of order \(a\ge2\), the vanishing denominator
factors are precisely those indexed by multiples of \(a\). Each has
a simple zero there. Thus the point is a pole exactly when \(a\le k\);
its order is then \(\lfloor k/a\rfloor\le\lfloor k/2\rfloor\).
By partial fractions, its coefficient contribution is a polynomial
of degree at most \(\lfloor k/2\rfloor-1\), multiplied by a number
of modulus one to the power \(v\).
Since
\[
\lfloor k/2\rfloor-1\le k-4=d-2
\qquad(k\ge5),
\]
the sum of these contributions is \(O_k((v+1)^{d-2})\).
For \(k=5\) the pole at \(-1\) has order two and contributes at most a linear error, exactly the allowed order.
Together with \eqref{eq:root-one-shift}, this proves
\eqref{eq:restricted-estimate} for \(v\ge0\).

For \(v<0\), the left side is zero.
The positive-part main term can be nonzero only in the bounded
interval \(-S/2<v<0\), where it is \(O_k(1)\).
For \(v\le-S/2\) both main quantities vanish.
This proves the uniform extension to all integers.
\end{proof}

\section{Uniform first differences}

Retain \(S=k(k+1)/2-1\), and set
\begin{equation}\label{eq:D}
D_{k,\ell}(n)=\coeff{n}(1-q)G_{k,\ell}(q).
\end{equation}
Thus \(D_{k,\ell}(n)\) is the difference between consecutive
coefficients of \(G_{k,\ell}\), with the coefficient at \(-1\) set
equal to zero.
Define
\begin{equation}\label{eq:shift-parameters}
L=\ell+\frac{k+1}{2},\qquad
\delta=\frac{k\ell+1}{2}-n,\qquad
x=\frac k2-\frac{\delta}{L}.
\end{equation}

Let \(B_k\) be the \(k\)-fold convolution of the indicator of
\([0,1]\), or equivalently the density of the sum of \(k\) independent
uniform random variables on that interval.
Inclusion--exclusion for a simplex intersected with the unit cube
gives
\begin{equation}\label{eq:spline}
B_k(x)=\frac1{(k-1)!}
\sum_{j=0}^{k}(-1)^j\binom{k}{j}(x-j)_+^{k-1}.
\end{equation}
For completeness, the volume of
\(\{u\in[0,1]^k:\sum u_i\le x\}\) is
\(\frac1{k!}\sum_j(-1)^j\binom{k}{j}(x-j)_+^k\);
differentiating gives \eqref{eq:spline}.

\begin{proposition}\label{prop:uniform}
For each fixed \(k\ge5\), uniformly over
\(0\le n\le\lfloor k\ell/2\rfloor\) as \(\ell\to\infty\),
\begin{equation}\label{eq:uniform}
D_{k,\ell}(n)=\frac{L^{k-2}}{k!}B_k'(x)
+O_k(L^{k-4}),
\end{equation}
where \(L,x\) are given by \eqref{eq:shift-parameters}.
\end{proposition}

\begin{proof}
Expanding the numerator in \eqref{eq:G} gives
\begin{equation}\label{eq:subset}
D_{k,\ell}(n)=
\sum_{J\subseteq[k]}(-1)^{|J|}
a_k\left(n-|J|\ell-\sum_{r\in J}r\right).
\end{equation}
Here \([k]=\{1,\ldots,k\}\).
All nonnegative arguments in the sum are at most
\(n=O_k(L)\). Lemma~\ref{lem:restricted} therefore introduces a total
error \(O_k(L^{k-4})\), since the number of subsets depends only on
\(k\).

For \(J\subseteq[k]\), put
\[
\varepsilon_J=\sum_{r\in J}\left(r-\frac{k+1}{2}\right).
\]
Using \(Lx=n+S/2\), direct calculation yields
\begin{equation}\label{eq:exact-shift}
n-|J|\ell-\sum_{r\in J}r+\frac S2
=L(x-|J|)-\varepsilon_J.
\end{equation}
For each \(j\in\{0,\ldots,k\}\),
\begin{equation}\label{eq:average-zero}
\sum_{|J|=j}\varepsilon_J=0.
\end{equation}
Indeed, for \(j\ge1\), each index occurs in
\(\binom{k-1}{j-1}\) subsets and
\(\sum_{r=1}^{k}(r-(k+1)/2)=0\).
The empty-subset case is immediate.

Let \(d=k-2\) and \(\phi(y)=y_+^d\).
Since \(d\ge3\), the function \(\phi\) is \(C^2\), including at zero.
Each \(\varepsilon_J\) is bounded in terms of \(k\), and Taylor's
theorem, uniformly for \(z=O_k(L)\), gives
\[
\phi(z-\varepsilon_J)
=\phi(z)-\varepsilon_J\phi'(z)+O_k(L^{d-2}).
\]
This follows from
\(\sup|\phi''(y)|=O_k(L^{d-2})\) on the relevant bounded-length
interval about \(z\), and remains valid if the interval crosses zero.
Take \(z=L(x-j)\) and sum over \(|J|=j\).
The linear terms vanish by \eqref{eq:average-zero}, leaving
\begin{equation}\label{eq:taylor-group}
\sum_{|J|=j}\phi(L(x-j)-\varepsilon_J)
=\binom{k}{j}L^d(x-j)_+^d+O_k(L^{d-2}).
\end{equation}
Equations \eqref{eq:subset}--\eqref{eq:taylor-group} now imply
\[
D_{k,\ell}(n)=
\frac{L^d}{k!\,d!}
\sum_{j=0}^{k}(-1)^j\binom{k}{j}(x-j)_+^d
+O_k(L^{d-2}).
\]
Differentiating \eqref{eq:spline} proves \eqref{eq:uniform}.
\end{proof}

\begin{lemma}\label{lem:curvature}
For each \(k\ge5\), there is a constant \(\alpha_k>0\) such that
\begin{equation}\label{eq:derivative-lower}
B_k'(x)\ge\alpha_k(k/2-x)
\qquad(1/2\le x\le k/2).
\end{equation}
\end{lemma}

\begin{proof}
The density \(B_2\) is triangular. Convolution gives, for \(j\ge3\),
\begin{equation}\label{eq:convolution}
B_j'(x)=B_{j-1}(x)-B_{j-1}(x-1).
\end{equation}
Each \(B_j\) is supported on \([0,j]\), is positive in its interior,
and is symmetric about \(j/2\).
We show by induction that \(B_j'(x)>0\) on \((0,j/2)\).
For \(x<1\), the second term in \eqref{eq:convolution} is zero and
the first is positive. For \(1\le x<j/2\), both arguments are in
the support, and \(x\) is nearer to \((j-1)/2\) than \(x-1\), because
\[
\left(x-\frac{j-1}{2}\right)^2
-\left(x-1-\frac{j-1}{2}\right)^2=2x-j<0.
\]
The symmetric strict unimodality of \(B_{j-1}\) gives the claim.
The triangular case supplies the induction base, and symmetry
gives the negative derivative on the right half.

For \(k\ge5\), differentiation of \eqref{eq:convolution} at the
center is valid and gives
\[
B_k''(k/2)=
B_{k-1}'(k/2)-B_{k-1}'(k/2-1)
=2B_{k-1}'(k/2)<0.
\]
Consequently \(B_k'(x)/(k/2-x)\) is positive on
\([1/2,k/2)\) and extends continuously to
\(-B_k''(k/2)>0\) at its right endpoint.
Its positive minimum on this compact interval proves
\eqref{eq:derivative-lower}.
\end{proof}

\begin{corollary}\label{cor:bulk}
For fixed \(k\ge5\), there are constants \(\gamma_k>0\) and
\(\ell_1(k)\) such that
\begin{equation}\label{eq:bulk-lower}
D_{k,\ell}(n)\ge\gamma_k L^{k-3}
\end{equation}
whenever \(\ell\ge\ell_1(k)\) and
\(\ell\le n\le\lfloor k\ell/2\rfloor\).
\end{corollary}

\begin{proof}
For these indices, \(\delta\ge1/2\).
Also \(x\ge(\ell+S/2)/L\ge1/2\), and \(x<k/2\).
Proposition~\ref{prop:uniform} and Lemma~\ref{lem:curvature} give
\[
D_{k,\ell}(n)
\ge\frac{\alpha_k}{k!}\delta L^{k-3}-C_kL^{k-4}
\ge\frac{\alpha_k}{2k!}L^{k-3}-C_kL^{k-4}.
\]
For sufficiently large \(L\), this is at least
\(\alpha_kL^{k-3}/(4k!)\).
\end{proof}

\begin{remark}
When \(k\ell\) is odd, the first difference between the two equal
middle coefficients has index \((k\ell+1)/2\), outside the range
\(n\le\lfloor k\ell/2\rfloor\) used here.
The half-unit offset in
\(\delta=(k\ell+1)/2-n\) matters at the middle coefficient.
A remainder only of order \(L^{k-3}\) would be insufficient there.
The shifted leading term in Lemma~\ref{lem:restricted} and the
cancellation in \eqref{eq:average-zero} give the smaller uniform
remainder needed in Corollary~\ref{cor:bulk}.
\end{remark}

\section{Comparison and proof of the conjecture}

\begin{lemma}\label{lem:small}
Fix \(k\ge5\).
For all nonnegative integers \(\ell,s,t\) satisfying
\eqref{eq:alignment}, and all \(0\le n\le\lfloor k\ell/2\rfloor\),
\begin{equation}\label{eq:small}
\left|\coeff{n}(1-q)q^sG_{k-2,t}(q)\right|
\le C_k(\ell+1)^{k-4}.
\end{equation}
\end{lemma}

\begin{proof}
Put \(h=k-2\ge3\).
For \(v\ge0\), a solution of
\[
2z_2+3z_3+\cdots+hz_h=v,\qquad z_j\in\NN,
\]
is determined by \(z_2,\ldots,z_{h-1}\), if \(z_h\) exists.
Therefore
\[
a_h(v)\le(v+1)^{h-2}.
\]
Expanding the smaller numerator gives
\[
\coeff{n}(1-q)q^sG_{h,t}(q)
=\sum_{J\subseteq[h]}(-1)^{|J|}
a_h\left(n-s-|J|t-\sum_{r\in J}r\right).
\]
Every nonnegative argument is at most \(n\le k\ell/2\).
The absolute value is thus at most
\[
2^h(k\ell/2+1)^{h-2}
=O_k((\ell+1)^{k-4}),
\]
uniformly in \(s,t\).
\end{proof}

\begin{lemma}\label{lem:edge}
Fix an integer \(k\ge5\), and let \(\ell,s,t\in\NN\) satisfy
\eqref{eq:alignment}, including \(s\ne1\).
Then the polynomial \(F\) in \eqref{eq:F} satisfies
\begin{equation}\label{eq:edge}
\coeff{n}(1-q)F(q)\ge0\qquad(0\le n\le\ell).
\end{equation}
\end{lemma}

\begin{proof}
The numerator of \(G_{k,\ell}\) in \eqref{eq:G} has no nonconstant
term of degree at most \(\ell\). Thus
\(D_{k,\ell}(n)=a_k(n)\) in this range.
If \(s>n\), the shifted term contributes zero.
Otherwise \(s\le n\le\ell\), and \eqref{eq:alignment} implies
\[
t=\frac{k\ell-2s}{k-2}\ge\ell.
\]
It follows that \(0\le n-s\le t\), so the smaller numerator also has
no nonconstant contribution at the relevant index. Hence
\begin{equation}\label{eq:edge-difference}
\coeff{n}(1-q)F(q)=a_k(n)-a_{k-2}(n-s).
\end{equation}

For \(s=0\), inclusion of the allowed part sizes gives nonnegativity.
For \(s\ge2\), choose \(u,v\in\NN\) with \(s=2u+3v\).
Append \(u\) parts of size 2 and \(v\) parts of size 3 to each
partition counted by \(a_{k-2}(n-s)\).
This is an injection into the partitions counted by \(a_k(n)\):
removing the prescribed added parts recovers the original partition.
The required part sizes are available because \(k-2\ge3\).
Since \(s\ne1\), these cases exhaust the possibilities.
\end{proof}

\begin{proof}[Proof of Theorem~\ref{thm:main}]
Let \(N=k\ell\), and write \(F(q)=\sum_{n=0}^{N}c_nq^n\),
padding by zeros when needed.
The degree of \(G_{k-2,t}\) is \(N-2s\), so both terms in
\eqref{eq:F} are symmetric about \(N/2\). Thus \(c_n=c_{N-n}\).

Lemma~\ref{lem:edge} gives \(c_n-c_{n-1}\ge0\) for
\(0\le n\le\ell\), with \(c_{-1}=0\).
For \(\ell\le n\le\lfloor N/2\rfloor\),
Corollary~\ref{cor:bulk} and Lemma~\ref{lem:small} give
\[
c_n-c_{n-1}
\ge\gamma_kL^{k-3}-C_k(\ell+1)^{k-4}>0
\]
for sufficiently large \(\ell\), with a threshold depending only
on \(k\).
Since \(c_0\) is zero if \(s=0\) and is one otherwise, all
coefficients in the left half are nonnegative and nondecreasing.
Symmetry proves nonnegativity and unimodality for the whole sequence.

Finally suppose \(s=1\). Alignment forces \(\ell\ge1\).
The first two coefficients of \(G_{k,\ell}\) are both one, and
the shifted polynomial subtracts one at degree one. Consequently
\(c_0=1\) and \(c_1=0\). Symmetry gives \(c_N=1\), with \(N\ge5\).
A sequence that drops at the first step and subsequently rises
cannot be unimodal. This also covers width one.
\end{proof}

\section{Formal Verification}

The accompanying Lean~4 project proves Theorem~\ref{thm:main} and
translates the conclusion back to the original parameters. Its public
theorem \texttt{SZ.original\_conjecture} quantifies over integers
\(k,m,b\), uses the exponent in \eqref{eq:original}, and tests
unimodality through the actual degree of the difference. The degree
can drop when the outer coefficients cancel. The project proves that
trimming the terminal zeros preserves the required property.
Its Gaussian polynomials satisfy the denominator-cleared identity
\eqref{eq:G}, and the coefficients of \(H_k\) count multiplicity
vectors for the allowed parts.

The formal proof gives a second derivation of
Lemma~\ref{lem:restricted}. At height five define
\[
 D(q)=\prod_{j=2}^{5}(1-q^j),\qquad
 P(q)=343-860q+739q^2-216q^3 .
\]
Taking fourth differences of the cubic gives
\(\sum_{n\ge0}(n+7)^3q^n=P(q)/(1-q)^4\).
The project supplies an explicit \(E(q)\in\mathbb Z[q]\) and checks,
as a polynomial identity,
\[
 (1-q^{60})^2\bigl(720(1-q)^4-P(q)D(q)\bigr)
       =E(q)D(q)(1-q)^4.
\]
Consequently \(720H_5(q)-\sum_{n\ge0}(n+7)^3q^n
= E(q)/(1-q^{60})^2\). Each coefficient on the right has absolute
value at most \((n+1)\sum_i |[q^i]E(q)|\). This establishes the
base estimate at every index.

For the induction, put \(K_k=k!(k-2)!\), \(c_k=S/2\) for the current
height, and \(R_k(v)=K_ka_k(v)-(v+c_k)_+^{k-2}\).
For \(j=k+1\), one has \(c_j=c_k+j/2\),
\(K_j=j(k-1)K_k\), and \(a_j(v)-a_j(v-j)=a_k(v)\).
At \(v\ge0\), set \(z=v+c_k\). Since \(c_k\ge j/2\), the
binomial theorem gives
\[
 (z+j/2)^{k-1}-(z-j/2)^{k-1}
      =j(k-1)z^{k-2}+O_k((1+v)^{k-4}).
\]
The constant and quadratic terms in the increment cancel.
It follows that \(R_j(v)-R_j(v-j)=O_k((1+v)^{k-4})\) by the
induction hypothesis. Iteration reaches a negative index in at
most \(v+1\) steps, where the residual is bounded in terms of \(j\).
Thus \(R_j(v)=O_j((1+v)^{j-4})\), as required.

For Proposition~\ref{prop:uniform}, the formal proof varies the
subset shift by a real parameter \(u\), replacing \(x-|J|\) by
\(x-|J|-u\varepsilon_J\) in each truncated power. Reflection
\(i\mapsto k+1-i\) preserves subset size and negates
\(\varepsilon_J\), so the derivative of the sum at zero vanishes.
Its second derivative is bounded uniformly on
\([0,k]\times[0,1]\) in \((x,u)\). Setting \(u=1/L\) gives
the error in \eqref{eq:uniform} after multiplication by \(L^{k-2}\).

The source repository is \url{https://github.com/felixfei-research/stanley-zanello}.
The project uses Lean~4.24.0 and Mathlib at revision
\texttt{f897ebcf72cd} (v4.24.0).
The transitive axiom report for the original conjecture lists
\texttt{propext}, \texttt{Classical.choice} and \texttt{Quot.sound}.
Running \texttt{bash verify.sh} in \texttt{anc/lean} builds the project
with the pinned dependencies and checks the output of \texttt{\#print axioms}
for the main theorem and thirteen supporting declarations.

\section{Computational checks}

An accompanying Python program independently checks the identities and parameter
conventions using exact integer and rational arithmetic.
The finite scan
\[
5\le k\le16,\qquad 1\le\ell\le160,
\qquad 0\le t\le\left\lfloor\frac{k\ell}{k-2}\right\rfloor
\]
with the parity condition imposed contains \(149\,252\) admissible
tuples with \(s\ne1\).
Of these tuples, \(245\) fail unimodality or nonnegativity, all with
\(\ell\le25\). The scan gives no threshold beyond the tested range. Gaussian coefficients are cross-checked using a
separate \(q\)-Pascal recurrence, and all failing tuples are retained
with an exact negative first-difference witness.
The ancillary verifier regenerates the full scan and its witnesses.

\section*{Author Contributions and AI Disclosure}
The author initiated and directed the research, established the
problem-selection criteria, designed the verification workflow, and
evaluated and integrated the results.

GPT-6 Pro (OpenAI, via ChatGPT) identified the conjecture, developed
the mathematical proof, and wrote the initial verification code. GPT-6.2 Astra (OpenAI, via Codex) audited the
proof, developed the Lean~4 formalization, established an alternative
algebraic proof of the restricted-partition estimate, and performed
computational verification. Both models contributed to manuscript
preparation. Additional model runs cross-checked the arguments and
reviewed the related literature.

The author assumes responsibility for the final manuscript and its
submission.

\begin{thebibliography}{9}

\bibitem{Alpoge14}
L. Alpoge,
\emph{Proof of a conjecture of Stanley--Zanello},
J. Combin. Theory Ser. A \textbf{125} (2014), 166--176.
\href{https://doi.org/10.1016/j.jcta.2014.03.004}{doi:10.1016/j.jcta.2014.03.004}.
\href{https://arxiv.org/abs/1308.2358}{arXiv:1308.2358}.

\bibitem{BGK}
M. Beck, I. M. Gessel, and T. Komatsu,
\emph{The polynomial part of a restricted partition function related
to the Frobenius problem},
Electron. J. Combin. \textbf{8} (2001), no.~1, Note~7.
\href{https://doi.org/10.37236/1592}{doi:10.37236/1592}.
\href{https://arxiv.org/abs/math/0306036}{arXiv:math/0306036}.

\bibitem{CJ}
W. Y. C. Chen and I. D. D. Jia,
\emph{Semi-invariants of binary forms pertaining to a unimodality
theorem of Reiner and Stanton},
Int. J. Math. \textbf{32} (2021), no.~12, 2140003.
\href{https://doi.org/10.1142/S0129167X21400036}{doi:10.1142/S0129167X21400036}.
\href{https://arxiv.org/abs/2011.01467}{arXiv:2011.01467}.

\bibitem{Kang}
M. Kang,
\emph{Strict unimodality of a Gaussian difference for even $n$},
preprint (2026).
\href{https://mathdb.com/abs/2611.00169v1}{MathDB:2611.00169v1}.

\bibitem{KUW23}
C. Koutschan, A. K. Uncu, and E. Wong,
\emph{A unified approach to unimodality of Gaussian polynomials},
Proceedings of the 2023 International Symposium on Symbolic and
Algebraic Computation (ISSAC '23), 434--442.
\href{https://doi.org/10.1145/3597066.3597113}{doi:10.1145/3597066.3597113}.
\href{https://arxiv.org/abs/2302.04067}{arXiv:2302.04067}.

\bibitem{MPP}
S. Melczer, G. Panova, and R. Pemantle,
\emph{Counting partitions inside a rectangle},
SIAM J. Discrete Math. \textbf{34} (2020), no.~4, 2388--2410.
\href{https://doi.org/10.1137/20M1315828}{doi:10.1137/20M1315828}.
\href{https://arxiv.org/abs/1805.08375}{arXiv:1805.08375}.

\bibitem{RS98}
V. Reiner and D. Stanton,
\emph{Unimodality of differences of specialized Schur functions},
J. Algebraic Combin. \textbf{7} (1998), no.~1, 91--107.
\href{https://doi.org/10.1023/A:1008623312887}{doi:10.1023/A:1008623312887}.

\bibitem{SZ16}
R. P. Stanley and F. Zanello,
\emph{Some asymptotic results on \(q\)-binomial coefficients},
Ann. Comb. \textbf{20} (2016), no.~3, 623--634.
\href{https://doi.org/10.1007/s00026-016-0319-8}{doi:10.1007/s00026-016-0319-8}.
\href{https://arxiv.org/abs/1503.06367}{arXiv:1503.06367}.

\bibitem{SZ20}
R. P. Stanley and F. Zanello,
\emph{A generalization of a 1998 unimodality conjecture of Reiner
and Stanton},
J. Combin. \textbf{11} (2020), no.~1, 111--126.
\href{https://doi.org/10.4310/JOC.2020.v11.n1.a5}{doi:10.4310/JOC.2020.v11.n1.a5}.
\href{https://arxiv.org/abs/1711.10033}{arXiv:1711.10033}.

\end{thebibliography}
\end{document}
